
The following 65 primary studies constitute the corpus synthesised in this review. Studies discussed individually in the text are additionally listed in the References.

Álvarez Ariza, J., \& Hernández Hernández, C. (2026). Investigating the
Impact of Education 4.0 and Digital Learning on Students' Learning
Outcomes in Engineering: A Four-Year Multiple-Case Study. \emph{Informatics},
\emph{13}(2), 18. \url{https://doi.org/10.3390/informatics13020018}

Baranov, V. N., Bezrukikh, A. I., Konstantinov, I. L., Rudnitsky, E. A.,
Solopeko, N. S., \& Baykovskiy, Y. V. (2022). Digital Twin Application in
Teaching Students Majoring in Metallurgical Engineering. \emph{Vysshee
Obrazovanie v Rossii}, \emph{31}(2), 135--148.
\url{https://doi.org/10.31992/0869-3617-2022-31-2-135-148}

Barletta, V. S., Caivano, D., Catalano, C., Gentile, A., Posa, D. P., \&
Piccinno, A. (2025). Cyber Digital Twin to improve Security Education:
The RedTeam Knife approach. \emph{CHItaly 2025 - Proceedings of the 16th
Biannual Conference of the Italian SIGCHI Chapter}, 68.
\url{https://doi.org/10.1145/3750069.3750102}

Boltsi, A., Kosmanos, D., Xenakis, A., Chatzimisios, P., \& Chaikalis, C.
(2025). Towards a Novel Digital Twin Framework Proposal Within the
Engineering Design Process for Future Engineers: An IoT Smart Building
Use Case. \emph{Sensors (14248220)}, \emph{25}(11), 3504.
\url{https://doi.org/10.3390/s25113504}

Bonavolonta, F., Verde, M. T., Schiano Lo Moriello, R., \& Liccardo, A.
(2026). Sustainability comparison between a traditional electronics
measurement laboratory and an immersive metaverse laboratory: An
extended life cycle assessment analysis. \emph{Acta IMEKO}, \emph{15}(1), 1--nov.
\url{https://doi.org/10.21014/actaimeko.v15i1.2199}

Chen, S., \& Chen, Y. (2025). Research on the virtual simulation system
of intelligent manufacturing production line teaching based on digital
twin. \emph{Proceedings of 2025 International Conference on Artificial
Intelligence and Smart Manufacturing, ICAISM 2025}, 204--209.
\url{https://doi.org/10.1145/3756423.3756455}

Cortes, D., Ramirez, J., Villagomez, L., Batres, R., Vasquez-Lopez, V.,
\& Molina, A. (2020). Digital Pyramid: An approach to relate industrial
automation and digital twin concepts. \emph{Proceedings - 2020 IEEE
International Conference on Engineering, Technology and Innovation,
ICE/ITMC 2020}, 9198643.
\url{https://doi.org/10.1109/ice/itmc49519.2020.9198643}

Cruz Salazar, L. A., Tovar, H. S., Prada Salamanca, E. A., Márquez
Juárez, J. C., \& Cruz, E. M. (2025). Adaptive Learning and
Interoperability in I4.0: A Case Study of Agent-based Digital Twin with
Asset Administration Shell. \emph{IFAC-PapersOnLine}, \emph{59}, 202--207.
\url{https://doi.org/10.1016/j.ifacol.2025.12.150}

Cuiffi, J. D., Wang, H., Heim, J., Anthony, B. W., Kim, S., \& Kim, D. D.
(2021). Factory 4.0 Toolkit for Smart Manufacturing Training. \emph{ASEE
Annual Conference and Exposition, Conference Proceedings}.
\url{https://www.scopus.com/pages/publications/85124517580?origin=resultslist}

Eriksson, K., Alsaleh, A., Behzad Far, S., \& Stjern, D. (2022). Applying
Digital Twin Technology in Higher Education: An Automation Line Case
Study. \emph{Advances in Transdisciplinary Engineering}, \emph{21}, 461--472.
\url{https://doi.org/10.3233/atde220165}

Federal University of Amazonas, D. of P. E., Melo Freires, V. de, Lucas,
A., Silva, R. H. P., State University of Amazonas, D. of E. E., \&
Almeida, E. (2025). Development of a didactic solution for teaching
concepts related to Digital Twins using Educational Robot.
\emph{International Journal of Industrial Engineering and Management},
\emph{16}(4), 444--458. \url{https://doi.org/10.24867/ijiem-398}

Fu, X., Wang, Y., \& Fan, Y. (2025). Computer technology-enhanced digital
twin-driven innovation and practice of talent training mode of new
engineering experimental and training. \emph{Proceedings of The 2nd
International Conference on Intelligent Education and Computer
Technology, IECT 2025}, 690--697.
\url{https://doi.org/10.1145/3764206.3764314}

Fuertes, J. J., Prada, M. Á., Rodríguez-Ossorio, J. R., González-Herbón,
R., Pérez, D., \& Domínguez, M. (2021). Environment for Education on
Industry 4.0. \emph{IEEE Access}, \emph{9}, 144395--144405.
\url{https://doi.org/10.1109/access.2021.3120517}

Gallego, J., \& Vallverdú, F. (2025). Low-Cost Digital-Twin VR
Teleoperation of a 3D-Printed Robotic Arm for STEM Education. \emph{2025 11th
International Conference on Virtual Reality, ICVR 2025}, 351--356.
\url{https://doi.org/10.1109/icvr66534.2025.11172534}

García, D. C., Correa, I. A., González, S. G., Sánchez, A. Á., \&
González, G. V. (2026). Digital Twin for Designing Logic Gates in
Minecraft Through Automated Circuit Verification and Real-Time
Simulation. \emph{Electronics (Switzerland)}, \emph{15}(3), 499.
\url{https://doi.org/10.3390/electronics15030499}

González-Medina, J. A., González-Suárez, P., \& Hernández-Castellano, P.
M. (2026). Virtual Remotes Teaching Laboratories in the Progress of a
Sustainable Education. In \emph{Safe and Sustainable Value Creation by
Design} (pp.~421--429). \url{https://doi.org/10.1007/978-3-032-21154-5_46}

Guo, M., Zhang, Z., Orduña, P., \& Hussein, R. (2025). RHL-Butterfly: A
Scalable IoT-Based Digital Twinning Platform for Embedded Systems and
Remote Laboratories. \emph{International Journal of Online and Biomedical
Engineering}, \emph{21}(4), abr--28.
\url{https://doi.org/10.3991/ijoe.v21i04.54297}

Häfner, P., Bergmann, V., Häfner, V., Michels, F. L., Grethler, M., \&
Karande, A. (2024). An Adaptive Learning Environment for Industry 4.0
Competencies Based on a Learning Factory and Its Immersive Digital Twin.
\emph{International Conference on Computer Supported Education, CSEDU -
Proceedings}, \emph{1}, 720--731. \url{https://doi.org/10.5220/0012757000003693}

Hagedorn, L., Riedelsheimer, T., \& Stark, R. (2023). PROJECT-BASED
LEARNING IN ENGINEERING EDUCATION -- DEVELOPING DIGITAL TWINS IN A CASE
STUDY. \emph{Proceedings of the Design Society}, \emph{3}, 2975--2984.
\url{https://doi.org/10.1017/pds.2023.298}

Heidarinejad, M., \& Srivastava, A. (2023). Implication of Developing
Digital Twins to Improve Students' Learning Experiences. \emph{ASEE Annual
Conference and Exposition, Conference Proceedings}.
\url{https://www.scopus.com/pages/publications/85172150147?origin=resultslist}

Hu, W., Cao, Z., Messaoudi, A. E., Li, P., \& Zeng, Z. (2025). A desktop
learning factory for smart and resilient manufacturing based on digital
twin and AI. \emph{11th 2025 International Conference on Control, Decision
and Information Technologies, CoDIT 2025}, 739--744.
\url{https://doi.org/10.1109/codit66093.2025.11321707}

Jiang, F., Jiang, F., Xu, Y., Zhuang, Y., Qian, J., Liu, H., Su, X.,
Chen, S., \& Chen, X. (2026). \emph{A six-dimensional digital twin model for
flexible manufacturing systems in engineering education}.
\url{https://doi.org/10.1007/s44443-026-00589-7}

Keaveney, S., Athanasopoulou, L., Siatras, V., Stavropoulos, P.,
Mourtzis, D., \& Dowling, D. P. (2021). Development and Implementation of
a Digital Manufacturing Demonstrator for Engineering Education.
\emph{Procedia CIRP}, \emph{104}, 1674--1679.
\url{https://doi.org/10.1016/j.procir.2021.11.282}

Khan, M. S., Fatima, A., \& Irwin, J. (2026). In Teaching Lean
Manufacturing with Digital Twin: A Modern Approach to Engineering
Education. \emph{Journal of Engineering Technology}, \emph{43}(1), 28--38.
\url{https://doi.org/10.18260/jet.43.1.3}

Kumarasinghe, E. M. H. K., Liyanawaduge, N. N., Kulatunga, A. K.,
Seneviratne, A. M. N. D. B., \& Darmawardana, M. (2025). Transforming
Computer Integrated Manufacturing Facility to Use as a Digital Twin
Enabled Learning Factory. \emph{Lecture Notes in Mechanical Engineering},
881--888. \url{https://doi.org/10.1007/978-3-031-93891-7_97}

Kwateng, F., \& Leong, W. Y. (2026). Enhancing engineering learning
through digital twin and IoT applications: Student engagement and
industry alignment in the Ghanaian context. \emph{Eurasia Journal of
Mathematics, Science and Technology Education}, \emph{22}(3), em2787.
\url{https://doi.org/10.29333/ejmste/17970}

Li, X. (2025). Exploration of An Intelligent Teaching System for the
Post-Lithium Battery Course Empowered by AI. \emph{Education Reform and
Development,} \emph{7}(11), 302--308.
\url{https://doi.org/10.26689/erd.v7i11.13103}

Liljaniemi, A., \& Paavilainen, H. (2020). Using Digital Twin Technology
in Engineering Education -- Course Concept to Explore Benefits and
Barriers. \emph{Open Engineering}, \emph{10}(1), 377--385.
\url{https://doi.org/10.1515/eng-2020-0040}

Liljaniemi, A., \& Paavilainen, H. (2025). Enhancing engineering
education through immersive environments-a study of the Holodeck VR
system in hydraulics and pneumatics. \emph{Cogent Education}, \emph{12}(1),
2530900. \url{https://doi.org/10.1080/2331186x.2025.2530900}

Lin, T.-C., Chiu, C.-N., Wang, P.-T., \& Fang, L.-D. (2025). \emph{VisFactory:
Adaptive Multimodal Digital Twin with Integrated Visual-Haptic-Auditory
Analytics for Industry 4.0 Engineering Education}. \emph{1}(1), 3.
\url{https://doi.org/10.3390/multimedia1010003}

Lu, M., \& Hu, Z. (2025). Digital Twin-Enhanced Programming Education: An
Empirical Study on Learning Engagement and Skill Acquisition.
\emph{Computers}, \emph{14}(8), 322. \url{https://doi.org/10.3390/computers14080322}

Machado, G. S., Salgado, T. R. M., Ayres, F. A. C., Bessa, I. V.,
Medeiros, R. L. P., \& Lucena, V. F. (2025). Automatic Correction System
for Learning Activities in Remote-Access Laboratories in the
Mechatronics Area. \emph{Applied Sciences (Switzerland)}, \emph{15}(5), 2574.
\url{https://doi.org/10.3390/app15052574}

Morris, J., \& Wagner, J. R. (2023). Application of Extracurricular
Course Teaching Product Lifecycle Management Concepts to Undergraduates.
\emph{ASEE Annual Conference and Exposition, Conference Proceedings}.
\url{https://www.scopus.com/pages/publications/85172155593?origin=resultslist}

Mourtzis, D., Panopoulos, N., \& Angelopoulos, J. (2022). A hybrid
teaching factory model towards personalized education 4.0.
\emph{International Journal of Computer Integrated Manufacturing}, \emph{36}(12),
1739--1759. \url{https://doi.org/10.1080/0951192x.2022.2145025}

Nikolaev, S., Gusev, M., Padalitsa, D., Mozhenkov, E., Mishin, S., \&
Uzhinsky, I. (2018). Implementation of ``digital twin'' concept for modern
project-based engineering education. \emph{IFIP Advances in Information and
Communication Technology}, \emph{540}, 193--203.
\url{https://doi.org/10.1007/978-3-030-01614-2_18}

Ning, G., Luo, H., Yin, W., \& Zhang, Y. (2024). Exploration of the
Application and Practice of Digital Twin Technology in Teaching Driven
by Smart City Construction. \emph{Sustainability (Switzerland)}, \emph{16}(23),
10312. \url{https://doi.org/10.3390/su162310312}

Norambuena, N., Ortega, J., Rivera, F. M.-L., Covarrubias, M., Rivera,
J. L. V., Ramírez, E., \& Ketterer, C. I. G. (2025). Integrating Digital
Twins of Engineering Labs into Multi-User Virtual Reality Environments.
\emph{Applied Sciences (Switzerland)}, \emph{15}(7), 3819.
\url{https://doi.org/10.3390/app15073819}

Ortiz, J. S., Andaluz, V. H., \& Carvajal, C. P. (2026). Real-Time
Digital Twin Architecture for Immersive Industrial Automation Training.
\emph{Sensors}, \emph{26}(7), 2023. \url{https://doi.org/10.3390/s26072023}

Ortiz, J. S., Quishpe, E. K., Sailema, G. X., \& Guamán, N. S. (2025).
Digital Twin-Based Active Learning for Industrial Process Control and
Supervision in Industry 4.0. \emph{Sensors (14248220)}, \emph{25}(7), 2076.
\url{https://doi.org/10.3390/s25072076}

Pang, D., Cui, S., \& Yang, G. (2022). Remote Laboratory as an
Educational Tool in Robotics Experimental Course. \emph{International Journal
of Emerging Technologies in Learning}, \emph{17}(21), 230--245.
\url{https://doi.org/10.3991/ijet.v17i21.33791}

Park, J., Ho, C. Y., \& Li, L. K. B. (2026). AR Tunnel: An
augmented-reality digital twin for immersive learning of wind tunnel
laboratories. \emph{Computers and Education: X Reality}, \emph{8}, 100139.
\url{https://doi.org/10.1016/j.cexr.2026.100139}

Posio, J., Maljamäki, P., Haavikko, M., Tepsa, T., \& Väätäjä, H. (2023).
EXPERIENCES AND LEARNING OUTCOMES OF USING VIRTUAL REALITY IN BUILDING
SERVICES ENGINEERING EDUCATION. \emph{SEFI 2023 - 51st Annual Conference of
the European Society for Engineering Education: Engineering Education
for Sustainability, Proceedings}, 1098--1108.
\url{https://doi.org/10.21427/kcvd-ph95}

Promyoo, R., Alai, S., \& El-Mounayri, H. (2019). \emph{Innovative digital
manufacturing curriculum for industry 4.0}. \emph{34}, 1043--1050.
\url{https://doi.org/10.1016/j.promfg.2019.06.092}

Riess, C., Walter, M. S. J., Tyroller, M., \& Nierlich, R. (2023). A
DIGITAL LEARNING ENVIRONMENT TWIN OF A LAB ON PROTOTYPING TO GIVE
ENGINEERING STUDENTS DIGITAL ACCESS 24/7. \emph{SEFI 2023 - 51st Annual
Conference of the European Society for Engineering Education:
Engineering Education for Sustainability, Proceedings}, 2815--2823.
\url{https://doi.org/10.21427/g96d-we79}

Rigó, L., Fabianová, J., Čabaníková, L., \& Palinský, J. (2026). Digital
Twin of a Material Handling System Based on a Physical Construction-Kit
Model for Educational Applications. \emph{Machines}, \emph{14}(4), 429.
\url{https://doi.org/10.3390/machines14040429}

S, H. (2026). \emph{A Study On Impact Of Remote, Twin-Based Labs On Student
Performance And Engagement}. \emph{8}(1).
\url{https://doi.org/10.36948/ijfmr.2026.v08i01.66381}

Sakkas, K., Ntagka, N. E., Spyridakis, M., Miltiadous, A., Glavas, E.,
Tzallas, A. T., \& Giannakeas, N. (2025). Multiplayer Virtual Labs for
Electronic Circuit Design: A Digital Twin-Based Learning Approach.
\emph{Electronics (Switzerland)}, \emph{14}(16), 3163.
\url{https://doi.org/10.3390/electronics14163163}

Shehadeh, A., Alshboul, O., Abuaddous, M., Nimer, H. F., \& Otoom, M.
(2026). Emerging WEFE technologies in engineering education: Integration
of optimization, explainable AI, digital twins, and virtual reality.
\emph{Array}, \emph{30}, 100756. \url{https://doi.org/10.1016/j.array.2026.100756}

Speicher, T., DeFranco, J., Stricker, C., Kumara, S., \& Bilén, S.
(2026). Digital twin technology for engineering education: An
experiential learning case study. \emph{Industry and Higher Education}.
\url{https://doi.org/10.1177/09504222261443589}

Szántó, N., Monek, G. D., \& Fischer, S. (2024). Development of an
Immersive, Digital Twin-Supported Smart Reconfigurable Educational
Platform for Manufacturing Training: A Proof of Concept. \emph{Journal of
Engineering Management and Systems Engineering}, \emph{3}(4), 199--209.
\url{https://doi.org/10.56578/jemse030402}

Tarng, W., Wu, Y.-J., Ye, L.-Y., Tang, C.-W., Lu, Y.-C., Wang, T.-L., \&
Li, C.-L. (2024). Application of Virtual Reality in Developing the
Digital Twin for an Integrated Robot Learning System. \emph{Electronics
(Switzerland)}, \emph{13}(14), 2848.
\url{https://doi.org/10.3390/electronics13142848}

Terkaj, W., Kleine, K., \& Kuts, V. (2024). Virtual labs for higher
education in industrial engineering. \emph{Proceedings of the Estonian
Academy of Sciences}, \emph{73}(2), 142--149.
\url{https://doi.org/10.3176/proc.2024.2.07}

Terkaj, W., Pessot, E., Kuts, V., Bondarenko, Y., Pizzagalli, S. L., \&
Kleine, K. (2024). A framework for the design and use of virtual labs in
digital engineering education. \emph{AIP Conference Proceedings}, \emph{2989},
1--jun. \url{https://doi.org/10.1063/5.0189669}

Tjahyadi, H., Prasetya, K., \& Murwantara, I. M. (2023). Digital Twin
Based Laboratory for Control Engineering Education. \emph{International
Journal of Information and Education Technology}, \emph{13}(4), 704--711.
\url{https://doi.org/10.18178/ijiet.2023.13.4.1856}

Toivonen, V., Lanz, M., Nylund, H., \& Nieminen, H. (2018). The FMS
Training Center - A versatile learning environment for engineering
education. \emph{Procedia Manufacturing}, \emph{23}, 135--140.
\url{https://doi.org/10.1016/j.promfg.2018.04.006}

Usman, M., Usmani, A., \& Jadidi, M. (2023). Re-envisioning Field
Surveying as an Immersive Gamification Experience. \emph{Interaction Design
and Architecture(s)}, (58), 72--89.
\url{https://doi.org/10.55612/s-5002-058-003}

Vásquez-López, V., García-León, A. A., \& López-Suárez, Ó. (2025).
Enhancing Engineering Education through Digital Twins and
Interdisciplinary Collaboration. \emph{11th International Conference on
Higher Education Advances (HEAd'25)}, 864--871.
\url{https://doi.org/10.4995/head25.2025.20074}

Veitia, M. S.-I., Cousquer, C., Garcia, R., Khaoulan, S., Hauquier, F.,
Gomez, C., Caqueret, V., Pommet, M., Gervais, M., Lagarde, N., Garnier,
G. M., Guiga, W., Havet, J. L., Chapet, C., Habay, N., Bertrand, L.,
Haustant, C., Auchet, A., Roecker, B., \& Texier, A. (2025). When Virtual
Reality Reinvents Training: The Cnam Bet with CAP'VR. \emph{EPiC Series in
Computing}, \emph{107}, 384--394. \url{https://doi.org/10.29007/jjhj}

Wang, J., Fan, P., \& Zhang, H. (2025). Innovative Exploration of the
``Four-Dimensional Linkage'' Teaching System of Network Security Based on
Interdisciplinary Research. \emph{Proceedings of 2025 International
Conference on Educational Technology and Artificial Intelligence, ETAIC
2025}, 128--133. \url{https://doi.org/10.1145/3766557.3766580}

Wang, T., Wang, J., \& Chen, J. (2025). \emph{Digital Empowerment Reform of
Wind Farm Curriculum Design Based on Virtual Reality Integration}.
\emph{8}(3), 125--133. \url{https://doi.org/10.23977/curtm.2025.080318}

Wu, T. (2025). \emph{Teaching Reform Practices of Quality and Engineering
Management in the Context of Curriculum-based Ideological and Political
Education}. \emph{9}(2), 39--44. \url{https://doi.org/10.23977/aetp.2025.090207}

Xinyi, Z., \& Liye, L. (2025). \emph{Teaching Reform and Reflections on
Industrial IoT Curriculum under the Dual Needs of Industry 4.0 and New
Engineering Construction}. \emph{8}(5), 129--136.
\url{https://doi.org/10.23977/curtm.2025.080518}

Yin, R. (2025). Research on the Design and Development of Intelligent
Manufacturing Production Platform Based on Digital Twin. \emph{Proceedings of
2025 3rd International Conference on Internet of Things and Cloud
Computing Technology, IoTCCT 2025}, 106--110.
\url{https://doi.org/10.1145/3776865.3776884}

Zhang, M.-Q., Xu, Y., He, Y.-L., Zhang, Y., \& Song, X.-L. (2025).
AI-Driven Reconstruction and Optimization of Teaching Model for Fault
Detection and Diagnosis in Industrial Safety Production. \emph{Proceedings of
The 2nd International Conference on Intelligent Education and Computer
Technology, IECT 2025}, 234--239.
\url{https://doi.org/10.1145/3764206.3764241}

Zhou, Z., Oveissi, F., \& Langrish, T. (2024). Applications of augmented
reality (AR) in chemical engineering education: Virtual laboratory work
demonstration to digital twin development. \emph{Computers and Chemical
Engineering}, \emph{188}, N.PAG--N.PAG.
\url{https://doi.org/10.1016/j.compchemeng.2024.108784}
